% q0046.tex — When Ceteris Isn't Paribus
\question{
  id        = {0046},
  title     = {When Ceteris Isn't Paribus},
  source    = {OBECON},
  year      = {2026},
  round     = {Seletiva IEO -- Phase C},
  language  = {English},
  topic     = {Micro},
  subtopic  = {Econometrics / Selection Bias},
  type      = {objective},
  statement = {%
    A researcher compares wages: training programme participants earn \$24/hr
    vs.\ non-participants at \$18/hr, concluding the programme adds \$6/hr.
    However, 60\% of participants hold high school diplomas compared to 20\%
    of non-participants, and a diploma independently adds \$5/hr to wages.
    What is the actual programme effect on wages?
  },
  choices   = {%
    \item \$2/hr.
    \item \$4/hr.
    \item \$8/hr.
    \item It is not possible to determine the true value.
  },
  answer    = {B},
  solution  = {%
    Controlling for diploma composition: among participants without a diploma
    (40\%), expected wage $= 24 - 0.6 \times 5 = \$21$/hr. Among
    non-participants without a diploma (80\%), expected wage
    $= 18 - 0.2 \times 5 = \$17$/hr. The pure training effect is
    $21 - 17 = \$4$/hr. The remaining \$2 of the raw \$6 gap is attributable
    to selection bias (diploma endowment). (B) is correct.
  },
}
